% 介绍正文（无章标题；由 main.tex 的 \chapter* 或 promo.tex 引入）
% 在《技术组简介》《各组别宣传文案》的基础上，补充职责边界与工程场景。
\begingroup
% 避免长分点只有首行或末行落在另一页；设置仅用于本篇前言。
\clubpenalty=10000
\widowpenalty=10000

\subsection*{技术组如何协作}
一台机器人从图纸走到赛场，需要结构能够承载、执行机构能够受控、传感器能够提供信息，软件还要把这些能力组织起来。技术组通常分为机械、电控与视觉等组别，分工围绕各自负责的系统展开。

\begin{itemize}
    \item \textbf{机械组：让结构满足运动与安装需求。}负责机构设计、零件加工与装配，兼顾强度、重量、空间和维护便利性。例如，相机装在哪里、视野是否被遮挡、机器人运动时支架是否明显振动，都需要机械组参与解决。SolidWorks 等建模工具用于表达和验证设计。
    \item \textbf{电控组：让执行机构按指令稳定工作。}负责嵌入式软件、运动控制与设备通信，电路设计、焊接和供电调试也常由该组承担。例如，读取电机反馈、控制底盘速度、反馈云台姿态，以及接收上层软件的指令。常用工具包括 STM32 开发工具与电子设计自动化（EDA）软件。
    \item \textbf{视觉组：把传感器数据变成可用的状态与任务信息。}参与工业相机、镜头、激光雷达和计算平台的选型与调试，开发目标识别、跟踪、定位、导航及决策相关软件。软件以 C++ 为主，Python 常用于数据处理、训练与实验。视觉组既要关心算法结果，也要关心设备能否稳定采集、程序能否在机器人上持续运行。
\end{itemize}

以相机接入为例，视觉组提出成像需求，机械组落实安装位置和支架，电控组协助确认供电、通信及姿态反馈；三组共同验证运动时的数据是否可用。接口应约定单位、坐标系、时间戳与更新频率，否则各模块接在一起仍可能出错。本教材也将帮助视觉组新生逐步建立系统协作意识。

\subsection*{视觉组负责哪些任务}
理解视觉组，可以先看它需要回答哪些问题。下面按自瞄、导航、决策和雷达四类业务任务介绍；它们彼此协作，具体人员与模块的分工可以交叉。

\begin{itemize}
    \item \textbf{自瞄：目标在哪里，运动状态如何？}从相机图像中识别装甲板等目标，估计其空间位置，并通过连续观测维持跟踪、估计运动状态，向控制系统提供目标信息。检测解决“这一帧看到了什么”，跟踪解决“前后帧是否为同一目标、状态如何变化”。交付结果应包含有效状态；目标丢失时，也要让下游知道当前信息已经不可用。
    \item \textbf{导航：机器人在哪里，怎样到达目的地？}利用激光雷达、惯性测量单元（IMU）等传感器进行定位，结合地图、路径规划和局部避障，让机器人向给定目标移动。同步定位与建图（SLAM）用于同时估计自身位置与构建地图，导航系统则还要处理行进过程中的障碍和失败情况。ROS2 与 Nav2 是常见的软件基础。导航需要反馈到达、受阻或失败等状态，供上层决定下一步。
    \item \textbf{决策：何时去哪里，执行什么任务？}结合比赛状态、自身状态、感知结果和任务执行反馈，选择或切换机器人当前的行为。以哨兵为例，决策可以安排巡逻、前往指定区域或等待条件满足，再把目标交给导航等模块执行。状态机与行为树（Behavior Tree）是组织行为的常见方法。决策负责选择任务，导航负责完成移动；遇到目标不可达、信息过期等情况，决策还要决定重试、改选任务还是等待。
    \item \textbf{雷达：场上有哪些目标，它们分布在哪里？}借助较开阔的观测视角，识别与跟踪场上机器人，把观测结果转换到场地坐标中，并按通信约定共享目标位置等信息，为队伍提供全局态势参考。这里的“雷达”指比赛中的雷达站任务，不等同于导航所用的激光雷达传感器。其重点包括覆盖范围、目标身份关联、遮挡后的重新识别，以及共享信息的准确性与时效性。
\end{itemize}

这些任务共享一些方法与工具：传统图像处理和深度学习用于提取信息；OpenCV 是计算机视觉库，PyTorch 常用于模型训练；ROS2 提供节点通信等基础，Nav2 组织导航功能。学习时先明确问题与输入输出，再选择工具。

\subsection*{日常工作：从采集数据到交付系统}
按以下环节梳理一个功能，可以帮助你判断问题出在数据、算法，还是系统集成上。

\begin{enumerate}
    \item \textbf{明确需求与评价方式。}先描述功能在什么场景下工作、接收什么数据、输出给谁。例如，识别任务需要统计漏检和误检，跟踪任务还要观察目标切换与丢失后的恢复，导航任务需要记录是否到达及失败原因。评价方式应在实验前确定，才能比较改动是否有效。
    \item \textbf{接入设备并建立可靠输入。}根据视野、距离、接口带宽、算力与安装条件选择传感器和计算平台；完成驱动配置、曝光调节、相机标定及安装关系测量。保存能代表实际场景的原始数据，并检查时间戳、丢帧和设备断连情况。后续算法的效果首先受输入质量约束。
    \item \textbf{实现算法并验证中间结果。}根据任务选择图像处理、学习模型或状态估计等方法，先用记录的数据验证。除了最终输出，还应检查检测框、坐标变换、跟踪状态或规划路径等中间结果。使用学习模型时，还需要整理和标注数据、划分训练与评估数据，避免用同一批样本同时训练和证明效果。
    \item \textbf{接入机器人并完成联调。}把算法封装为可调用的模块或 ROS2 节点，接入设备、通信与控制链路，核对字段、单位、坐标和频率。补全启动配置、参数文件与异常处理，在真实计算平台上检查运行速度、资源占用和连续运行情况。
    \item \textbf{复盘问题并留下可复现的交付物。}保留代码版本、参数、设备配置、测试数据和结论，提供启动与排错说明。发现问题后，尽量形成可以回放或重复执行的案例，再验证修改。交付的目标是让队友能够运行、理解和继续维护这项功能。
\end{enumerate}

实际研发会在这些环节之间反复迭代。例如，远处目标总是漏检，既可能需要补充训练样本，也可能是镜头视野过大、目标像素太少。先查看原始图像和中间结果，才能判断应该改哪一层。

\subsection*{能力如何逐步建立}
入组时无需全部掌握以下能力：先学会运行与修改小程序，再进入专项，在项目中补齐知识。

\begin{itemize}
    \item \textbf{第一步：能够独立完成基础开发。}熟悉 Linux 下的文件、终端与进程操作，理解 C++ 的基本语法、类、容器和资源管理，能够使用 CMake 构建程序、借助 Git 管理改动。进一步学习智能指针与多线程，理解对象生命周期和共享数据问题。这些是卷一的主要训练内容。
    \item \textbf{第二步：能够解释数据与算法的含义。}逐步补充向量、矩阵、坐标变换与概率统计基础，理解图像、相机模型和常见视觉处理流程，再根据任务学习 OpenCV、深度学习、状态估计等知识。卷二提供相关基础与专项内容；学习时应能说清一个量表示什么、属于哪个坐标系、结果怎样验证。
    \item \textbf{第三步：能够负责一段完整功能。}结合自瞄、导航、决策或雷达任务，学习所需的 ROS2、规划、行为组织等技术，把输入、处理、输出和异常反馈接起来。卷三帮助理解赛场任务与工程协作，再将所学用于联调和复盘。
    \item \textbf{贯穿始终：能够把问题讲清楚。}记录复现步骤，阅读日志，使用调试器和可视化工具，并在求助时给出环境、现象、预期和已尝试的方法。对队友说明接口变化，对实验保留证据。这些习惯从第一份练习就可以开始培养。
\end{itemize}

\subsection*{主要挑战：为什么演示成功还不够}
单张图片上识别正确、一次路线顺利跑通，只验证了有限条件下的功能。机器人在实际运行中还会受到环境变化、运动、通信与资源限制的共同影响。下面列出几类需要持续验证的问题。

\begin{enumerate}
    \item \textbf{环境变化会破坏原本有效的观测。}
    光照改变、反光、运动模糊、远距离小目标和遮挡，都可能造成漏检、误检或跟踪中断。问题会继续传到下游，使目标状态跳变或失效。调试时应先检查原始图像，再按场景收集失败案例，分别判断是成像质量、检测还是跟踪环节的问题。修改后要回放多类场景，检查改善某一场景时是否损害了其他场景。

    \item \textbf{数据“看起来合理”，仍可能没有对齐。}
    相机、激光雷达与机器人本体各有坐标系，安装关系和标定参数会影响空间转换；不同设备的数据还可能对应不同时间。静止时误差不明显，运动后却可能表现为目标位置漂移或地图重影。需要核对单位、坐标轴方向、变换关系与时间戳，并在静止和运动条件下分别验证。安装位置变化后，也应重新检查有关参数。

    \item \textbf{单个算法运行快，不代表整条链路响应及时。}
    采集、传输、处理、排队和发送都会消耗时间。程序即使持续输出，也可能正在处理积压的旧数据；偶发卡顿还会使输出间隔忽长忽短，影响下游使用。应记录各环节耗时、数据采集到输出的时间差，以及延迟波动，定位瓶颈后再优化。队列如何处理过期数据、下游如何识别过期结果，都需要明确约定。

    \item \textbf{离线评估与实际场景可能存在差距。}
    训练或调试数据可能只覆盖固定场地、视角与光照，实机却会遇到新的距离、背景和遮挡组合。评估数据若与训练数据过于相似，也容易高估效果。应将独立采集的场景用于评估，并持续补充实机失败样本，分别统计不同条件下的表现。每次调整都要在保留的测试数据上比较，避免只凭少数成功画面判断效果。

    \item \textbf{计算平台与设备需要经受持续运行。}
    开发电脑上的程序迁移到机器人后，会受到算力、内存、功耗和散热条件的限制；多个模块同时运行，还可能争用资源。设备重连、日志持续写入和长时间运行也会暴露短时测试没有发现的问题。应在目标平台上同时运行实际所需模块，观察资源占用、温度、错误日志与数据连续性，并验证重新启动后的设备连接和配置是否正确。

    \item \textbf{任务受阻时，系统仍要有明确行为。}
    路径被挡住、定位失效、目标消失或状态更新中断，都可能让正常流程无法继续。如果只设计成功路径，机器人可能反复切换任务、持续等待，或继续使用过期信息。应为任务定义开始、完成、失败和超时条件，规定重试、取消或等待的处理方式。联调时主动构造受阻与数据中断场景，检查各模块是否给出一致的状态反馈。

    \item \textbf{跨组问题需要共同的证据来定位。}
    相似现象可能来自不同原因：画面抖动可能与支架振动有关，输出异常可能来自参数、通信或软件版本不一致。应记录测试时的代码版本、硬件状态、配置与日志，沿数据流逐段检查。接口变更要同步更新文档并通知相关成员；修复后保留复现案例，方便验证同类问题是否再次出现，也让后续成员能够接手。
\end{enumerate}

通识阶段先练好一件事：让程序稳定运行，能够解释输入输出，遇到错误时留下可复现的记录。再逐步扩大负责的范围，把所学与机器人上的实际问题连接起来。
\par
\endgroup
